\documentclass[11pt,letterpaper]{article}
\usepackage[T1]{fontenc}
\usepackage{lmodern,microtype,amsmath,amssymb,amsthm,mathtools,booktabs,xcolor}
\usepackage[margin=1.02in,top=0.94in,bottom=0.97in]{geometry}
\usepackage[hidelinks]{hyperref}
\usepackage{fancyhdr}
\usepackage{enumitem}
\hypersetup{pdftitle={How often can two products of Dirichlet L-functions agree locally?},pdfauthor={OpenAI ChatGPT}}
\newcommand{\Fp}{\mathbf F_p}
\newcommand{\Z}{\mathbf Z}
\newtheorem{theorem}{Theorem}[section]
\newtheorem{lemma}[theorem]{Lemma}
\newtheorem{corollary}[theorem]{Corollary}
\newtheorem{proposition}[theorem]{Proposition}
\theoremstyle{remark}
\newtheorem{remark}[theorem]{Remark}
\definecolor{ink}{HTML}{173A48}
\definecolor{ruleblue}{HTML}{91ADB5}
\pagestyle{fancy}\fancyhf{}
\fancyhead[L]{\small\color{ink}Local Euler factors of Dirichlet products}
\fancyhead[R]{\small\color{ink}Matching densities}
\fancyfoot[C]{\small\thepage}
\renewcommand{\headrulewidth}{0.3pt}
\setlength{\headheight}{14pt}
\setlength{\emergencystretch}{1.8em}
\raggedbottom
\begin{document}
\thispagestyle{plain}
\noindent{\small\sffamily\color{ink}ANALYTIC NUMBER THEORY AND FINITE GEOMETRY\hfill 27 NOVEMBER 2026}\par
\vspace{.35em}\noindent{\color{ruleblue}\rule{\textwidth}{.6pt}}\par
\vspace{.85em}
\noindent{\huge\bfseries\color{ink}How often can two products of\par
\noindent Dirichlet \(L\)-functions agree locally?}\par
\vspace{.6em}
\noindent{\large Exact matching densities for cubic and quintic characters}\par
\vspace{1.25em}
\begin{abstract}\noindent
Two different products of fifteen quintic Dirichlet $L$-functions can have
identical Euler factors at $113/125$ of the primes. This is best possible
in a character space of rank three. More generally, we compute the maximum
agreement density for every packet size and ambient rank at character
orders three and five. For arbitrary odd prime order $p$, the leading
base-$p$ digit $b$ of the packet size gives a lower bound on disagreement.
Before saturation, equality is possible exactly when $b$ divides $p-1$
or $p+1$. We describe the extremal packets, prove a gap when the divisibility
test fails at small sizes, and give a complete two-factor spectrum.
For factors of one primitive conductor the sharp large-size bound increases by a factor $2^{r-1}$. For cubic characters the exact primitive minimum at every size is $(2/3)^{\min(r,\lfloor\log_2(2k)\rfloor)}$. The density results are unconditional. The Dirichlet GRH theorem in
ProveIt gives a finite distinguishing prime and, for extremal pairs, a
weighted prime count with an exact Fourier norm.
\end{abstract}

\section{The question and the input}

At a prime $\ell$ away from the modulus, a product
$\prod_{i=1}^k L(s,\chi_i)$ records the multiset
$\{\chi_1(\ell),\ldots,\chi_k(\ell)\}$ in its local Euler polynomial.
Agreement forgets which factor supplied which value. The question is how
often two different lists can be indistinguishable after this loss of labels.
Here $r$ is the rank of the ambient character space. A packet may occupy a
proper subspace, and characters may repeat.

The merged ProveIt article solves the quadratic extremal problem and describes
part of its high-agreement spectrum. It explicitly asks for the corresponding
odd-prime densities in the closing questions of its recovery section
\cite{Article}. Theorem~\ref{thm:digit} answers the extremal question
for cubic and quintic characters for every $k,r$.
Theorem~\ref{thm:descent} identifies precisely which leading digits give
equality for an arbitrary odd prime. These theorems concern the maximum
possible agreement, with the whole set of possible densities given
separately for packets of size two in Proposition~\ref{prop:two}. Theorems~\ref{thm:primitive} and \ref{thm:cubic} give the sharp common-conductor results for primitive factors.

\paragraph{What the Lean source supplies.}
At ProveIt commit \texttt{2a9f8e3ee7}, the vendored readable
\texttt{grh} directory and upstream commit \texttt{72fead102f}
have the same Git tree. Its \texttt{Statements.lean} imports Mathlib.
The \texttt{Main.lean} exports include Mathlib's RH, right-half-plane
nonvanishing for Dirichlet characters with the principal pole exception,
and critical-line location of the open-strip zeros for every positive
modulus, including imprimitive characters \cite{Lean}. The proof obtains
a one-sided bound for a filtered prime sum and converts it to nonvanishing.
We use its zero-location conclusion in Section~\ref{sec:primes}, where the
smoothed explicit formula needs only Dirichlet characters. The upstream
README reports a Lean 4.35.0 build and independent kernel checks. We
inspected the pinned statements and tree and relied on the reported
verification. The new results below have written proofs and finite
computational checks.

\paragraph{Relation to previous mathematics.}
The planar bound for packets smaller than $p$ follows already from Lev's
bound on perfect directions for signed weights \cite{Lev}. His examples
include rotational configurations. Blokhuis, Korchm\'aros and Mazzocca
classify the relevant configurations with a collinear component
\cite{BKM}. The polynomial weight and shape theorems used below are
results about projective Reed--Muller codes \cite{GL}. Chen uses a related
Newton-sum argument for cyclotomic lattices \cite{Chen}. General Fourier
uncertainty inequalities, such as Meshulam's \cite{Meshulam}, constrain
support size. The packet problem also constrains total positive integral
mass, which can be much larger than support size. The contribution here is
to relate that mass to cyclotomic divisibility in every direction, descend
an equality case to the plane, and convert the result back to prime
densities. We found no statement of this all-rank packet formula in the
sources examined.
\section{The extremal density}

Let $p$ be an odd prime, let $X$ be an $r$-dimensional $\Fp$-space of Dirichlet characters of exponent dividing $p$ modulo $Q$, and suppose $r\ge1$. Write $V=X^*$, fix $\zeta=\exp(2\pi i/p)$, and identify evaluation at a reduced residue class $u$ with $v\in V$ by $\chi(u)=\zeta^{v(\chi)}$. The evaluation map onto $V$ is surjective. Indeed, its image is a subspace and its annihilator consists only of the trivial character. Each $v$ therefore receives equally many reduced residue classes. The prime number theorem in arithmetic progressions turns any proportion of $V$ into the natural density of the corresponding primes.

A packet is a multiset of $k$ characters from $X$. Two packets $A,C$ disagree at $v$ when their local polynomials $\prod_{x\in A}(1-\zeta^{v(x)}T)$ and $\prod_{x\in C}(1-\zeta^{v(x)}T)$ differ. Write $\delta(A,C)$ for the proportion of such $v$, equivalently the density of disagreement among unramified primes. Common entries may occur in both packets.

\begin{theorem}[A bound from the leading base-$p$ digit]\label{thm:digit}
Let $A\ne C$ have size $k\ge1$, and define
\[
 a=\lfloor\log_p k\rfloor,\qquad b=\lfloor k/p^a\rfloor .
\]
Then
\[
 \delta(A,C)\ \ge\ D_p(r,k):=
 \begin{cases}
 (p-1)(p+1-b)/p^{a+2},& a\le r-2,\\[2pt]
 (p-1)/p^r,&a\ge r-1.
 \end{cases}
\]
In the second range this is the exact minimum. In the first range it is the exact minimum whenever $b\mid(p-1)$ or $b\mid(p+1)$. Consequently $D_p(r,k)$ is the exact minimum for every rank and size at $p=3$ and $p=5$.
\end{theorem}

For the proof, let $c_x$ be the multiplicity of $x$ in $A$ minus its multiplicity in $C$. Thus $\sum c_x=0$, $\sum |c_x|\le 2k$, and $c\ne0$. Consider
\[
 H(v)=\sum_{x\in X}c_x\zeta^{v(x)} .
\]
For $h_j(v)=\sum_{x:\,v(x)=j}c_x$, equality of local polynomials means that every $h_j(v)$ vanishes. If $H(v)=0$, the integer polynomial $\sum_{j=0}^{p-1}h_j(v)T^j$ is a scalar multiple of $1+T+\cdots+T^{p-1}$, and its value at $1$ is zero. Hence $H(v)=0$ is equivalent to local agreement. Fourier inversion shows that $H$ is nonzero somewhere. Also $H(0)=0$, and the automorphism $\zeta\mapsto\zeta^u$ carries $H(v)$ to $H(uv)$ for $u\in\Fp^\times$. In particular the disagreement set is a union of punctured lines through the origin.

\subsection{An arithmetic bound on the first nonzero reduction}

Set $\mathcal O=\Z[\zeta]$ and $\pi=\zeta-1$. We use $\mathcal O/(\pi)=\Fp$ and $(p)=(\pi)^{p-1}$. The second identity follows from $p=\prod_{j=1}^{p-1}(1-\zeta^j)$, since each quotient $(1-\zeta^j)/(1-\zeta)$ is a unit, as also follows by using the inverse of $j$ modulo $p$. All $H(v)$ are divisible by $\pi$ since $\sum c_x=0$. Let $t$ be the smallest $\pi$-adic valuation among their nonzero values.

\begin{lemma}\label{lem:massdigit}
Write $t=j(p-1)+s$, where $j\ge0$ and $1\le s\le p-1$. Then $k\ge p^j s$. In particular $t\le a(p-1)+b$.
\end{lemma}

\begin{proof}
Take $v$ attaining $t$ and abbreviate its vector of histogram differences by $h$. Evaluation at $\zeta$ gives an additive isomorphism
\[
 \{(n_0,\ldots,n_{p-1})\in\Z^p:\textstyle\sum_i n_i=0\}
 \longrightarrow \pi\mathcal O,\qquad n\longmapsto\sum_i n_i\zeta^i.
\]
The cyclotomic polynomial calculation above proves injectivity. The differences $\zeta^{i+1}-\zeta^i$ for $0\le i\le p-2$ give a $\Z$-basis of $\pi\mathcal O$, proving surjectivity. Now $H(v)\in p^j\pi\mathcal O$, so $h=p^j n$ for an integral zero-sum vector $n$. Its positive and negative parts each have mass $m\le k/p^j$. The number $\sum_i n_i\zeta^i$ has valuation $s$.

Expand each $\zeta^i=(1+\pi)^i$. For $0\le \ell<s$ one obtains
\[
  \sum_i n_i\binom{i}{\ell}\equiv0\pmod p .
\]
Indeed, the first coefficient of index less than $s$ that survived modulo $p$ would give the unique term of smallest valuation. All earlier coefficients are multiples of $p$ and have valuation at least $p-1\ge s$. The binomial polynomials of degrees up to $s-1$ span the polynomials of those degrees over $\Fp$. Thus the two $m$-element multisets represented by the positive and negative parts of $n$ have identical power sums through degree $s-1$. If $m<s$, Newton's identities through order $m<p$ give identical elementary symmetric polynomials. Equality of the resulting monic polynomials forces equality of the multisets, contradicting $H(v)\ne0$. It follows that $k\ge p^jm\ge p^js$. The numbers $p^js$ are strictly increasing in the ordering $j(p-1)+s$, proving the final assertion.
\end{proof}

\begin{lemma}\label{lem:reductiondegree}
The function $f(v)=H(v)/\pi^t\pmod\pi$ is a nonzero polynomial function on $V$ of total degree at most $t$. It vanishes at the origin and its support is a union of punctured lines.
\end{lemma}

\begin{proof}
Choose dual coordinates and lift each $v_i$ to $\{0,\ldots,p-1\}$. In the product expansion
\[
  \zeta^{v(x)}
  =\prod_{i=1}^r\sum_{u_i=0}^{p-1}\binom{v_i}{u_i}
     (\zeta^{x_i}-1)^{u_i},
\]
the coefficient of the tensor binomial indexed by $u$ is divisible by $\pi^{|u|}$, where $|u|=\sum_i u_i$. This also holds after summing against $c_x$. Each coefficient of $H$ in that basis is an iterated forward difference at the origin, hence an integral combination of the values of $H$ on the grid. All these coefficients are therefore divisible by $\pi^t$ as well. Dividing by $\pi^t$ and reducing kills the terms with $|u|>t$, and $u_i!$ is invertible modulo $p$. A point attaining valuation $t$ gives $f\ne0$. Galois conjugation preserves $\pi$-adic valuation, so $f(v)\ne0$ if and only if $f(uv)\ne0$ for each $u\in\Fp^\times$.
\end{proof}

\subsection{Counting the surviving directions}

For $m,D\ge0$, put
\[
  W(m,D)=\min_{\substack{0\le e_i\le p-1\\ \sum_i e_i\le D}}
               \prod_{i=1}^m (p-e_i),
\]
with $W(0,D)=1$. If $D\ge m(p-1)$ then $W(m,D)=1$. Otherwise write $D=j(p-1)+h$, $0\le h<p-1$. In that case
\[
    W(m,D)=(p-h)p^{m-j-1}.
\]
To see this, transfer one unit from the smaller to the larger of two positive, unsaturated exponents. The degree stays fixed and the product decreases. Thus a minimum fills coordinates successively up to $p-1$.

\begin{lemma}\label{lem:projectivecount}
Suppose a nonzero polynomial function on $\Fp^{m+1}$ of degree at most $t\ge1$ vanishes at the origin and has support a union of punctured lines. Its support contains at least $W(m,t-1)$ projective directions.
\end{lemma}

\begin{proof}
A nonzero polynomial function of degree at most $D$ in $m$ affine variables is supported on at least $W(m,D)$ points. Reduce coordinate exponents to at most $p-1$ using $X_i^p=X_i$. If its highest exponent in the first variable is $e_1$, above each point where the leading coefficient is nonzero at least $p-e_1$ values survive. Apply the same argument to that coefficient. The exponents so selected form a monomial of degree at most $D$, yielding the claimed product lower bound.

We induct on the projective dimension $m$, with the case $m=0$ immediate. If some projective hyperplane misses the support, take it to be $x_0=0$. For a representing polynomial $F$, the term $F(0,x')$ vanishes as a function at every $x'$. Subtracting it writes the same function as $x_0G(x_0,x')$, with $\deg G\le t-1$. The affine chart $x_0=1$ contains one point from every supported line, so the affine bound gives the claim. If every hyperplane meets the support, induction within each hyperplane gives at least $W(m-1,t-1)$ supported directions there. Double counting incidences gives at least
\[
    \frac{p^{m+1}-1}{p^m-1}\,W(m-1,t-1)
    >p\,W(m-1,t-1)\ge W(m,t-1)
\]
directions overall. The last inequality follows by adjoining an exponent zero in the defining minimum.
\end{proof}

This count agrees with the minimum-distance formula for projective Reed--Muller codes in the homogeneous case \cite{GL}. It also applies directly to the reduction in Lemma~\ref{lem:reductiondegree}.

\begin{proof}[Proof of the lower bound in Theorem~\ref{thm:digit}]
The support of $H$ contains that of $f$. Lemmas~\ref{lem:massdigit}--\ref{lem:projectivecount} give at least $W(r-1,a(p-1)+b-1)$ projective directions, each containing $p-1$ nonzero $v$. The displayed formula for $W$ gives $D_p(r,k)$ after division by $p^r$.
\end{proof}

\subsection{Packets attaining the bound}

We first build packets of size $b$ in rank two, with precisely $b$ of the $p+1$ projective directions of agreement.

\emph{Split construction.} If $b\mid(p-1)$, let $H_0\subset\Fp^\times$ have order $b$. Take the packets $H_0 e_1$ and $H_0 e_2$ on independent coordinate axes. For a nonzero $v=(u,w)$ the two projected multisets agree exactly when both entries are nonzero and $u/w\in H_0$. This gives $b$ directions.

\emph{Nonsplit construction.} If $b\mid(p+1)$, identify the rank-two space with $K=\mathbf F_{p^2}$ and its dual via $(v,x)\mapsto\operatorname{Tr}_{K/\Fp}(vx)$. Let $U$ be the subgroup of norm-one elements, $H_0\le U$ its subgroup of order $b$, and $c\in U\setminus H_0$. Use the packets $H_0$ and $cH_0$. For $v\ne0$, all elements of $vH_0$ and $vcH_0$ have the same norm. If the traces match as multisets, the trace of $v$ matches that of an element of $vcH_0$. Two elements with equal trace and norm are equal or Frobenius conjugate, since they solve the same monic quadratic. The cosets are disjoint, so $v^p\in vcH_0$, or $v^{p-1}\in cH_0$. Conversely, this condition implies $(vH_0)^p=vcH_0$, so all traces match with multiplicity. The map from $K^\times/\Fp^\times$ to $U$ given by $v\mapsto v^{p-1}$ is bijective. Again exactly $b$ projective directions give agreement.

\begin{proof}[Proof of sharpness in Theorem~\ref{thm:digit}]
When $a\le r-2$, choose a subspace $L$ of dimension $a$ independent of the rank-two construction. Replace each character $x$ of either packet by all characters $x+\ell$ for $\ell\in L$. The new Fourier difference is the old one times $p^a\mathbf 1_{L^\perp}(v)$, by character orthogonality. The two original coordinates remain uniform on $L^\perp$. Hence the disagreement proportion becomes $p^{-a}(p-1)(p+1-b)/p^2$. The packets have size $bp^a\le k$; pad both equally to size $k$.

For $a\ge r-1$, take a hyperplane $L\subset X$ and any $x\notin L$. The packets $L$ and $x+L$ disagree exactly on the punctured line $L^\perp\setminus\{0\}$ and each has size $p^{r-1}\le k$. This includes $r=1$ with $L=\{0\}$. Finally the allowed leading digits for $p=3$ are $1,2$, and those for $p=5$ are $1,2,3,4$. Each divides $p-1$ or $p+1$ for its corresponding prime.
\end{proof}
\section{When a small packet attains the bound}
\label{sec:geometry}

There is a geometric obstruction to attaining the density bound.  It is
already visible before the first power of the character order.  We work in
the exponent space $X=\Fp^r$ of the characters.  For a multiset $A$ in $X$
and a linear form $g\in V=X^*$, the projected multiset consists of the
elements $g(x)$ with $x\in A$, counted with multiplicity.  Write $E$ for
the set of $g$ on which the projected multisets of two packets agree.

\begin{theorem}[Rigidity below the character order]\label{base:rigid}
Let $r\ge 2$ and $1\le b<p$.  Two distinct packets of size at most
$b$ and of equal size have disagreement density at least
\[
 d_b=\frac{(p-1)(p+1-b)}{p^2}.
\]
Equality is attainable precisely when $b\mid p-1$ or $b\mid p+1$.
For $b\ge2$, every equality case consists, after cancellation, of two
disjoint $b$-element sets in a common affine plane.  The lines joining
opposite sets have exactly $b$ directions. Together with the corresponding
points at infinity the sets form a dual $3$-net of order $b$.
\end{theorem}

\begin{proof}
Cancel common entries, and call the remaining multisets $A,B$. Fix
$x\in A$.  For each $g\in E$ some $y\in B$ has $g(y)=g(x)$.  Consequently
\[
 E\subseteq\bigcup_{y\in B}(x-y)^\perp.
\]
A union of $m$ distinct central hyperplanes has at most
$p^{r-1}+(m-1)(p^{r-1}-p^{r-2})$ elements. Indeed each additional
hyperplane meets the first one in $p^{r-2}$ elements.  Taking $m\le b$
proves the density bound.  For $b=1$ any two distinct singleton packets
attain it.  Assume $b\ge2$ and equality holds.

For every $x$ all $b$ hyperplanes are distinct and their union equals $E$.
Thus $B$ is a set of size $b$, and interchanging $A,B$ gives the same
conclusion for $A$.  The hyperplanes pass through a common
$(r-2)$-space $K$.  For $b=2$ take their intersection.  For $b\ge3$,
equality forces each later hyperplane to intersect the entire previous
union in exactly $p^{r-2}$ elements. Its intersections with the first
two hyperplanes must coincide, hence it contains their intersection.
The hyperplanes in this pencil are determined by $E$: any different
hyperplane intersects the union in at most $bp^{r-2}<p^{r-1}$ points.
In particular the pencil is independent of $x$.

Every $x-y$, with $x\in A$ and $y\in B$, thus belongs to one of $b$ fixed
one-dimensional subspaces of the plane $K^\perp\subset X$.  Fixing a
point of each set shows that all the points belong to one affine translate
of this plane.  From every point of $A$ there is exactly one line of
each of the $b$ directions to $B$, and the same assertion holds with the
sets interchanged.  In the projective completion take $C$ to be the $b$
directions at infinity. Every line joining points of different sets
among $A,B,C$ contains exactly one point of each. This is the defining
incidence property of a dual $3$-net.

Theorem 5.1 of Blokhuis, Korchm\'aros and Mazzocca and the classification
in its proof apply since $b<p$ and $C$ is collinear \cite{BKM}.  They
place $A\cup B$ on a possibly reducible conic.  In the two-line case
the directions are differences of parameters on parallel lines, or
ratios on intersecting lines.  On a nonsingular conic the corresponding
group of parameters is the additive group of $\Fp$, the multiplicative
group of $\Fp$, or the norm-one group of $\mathbb F_{p^2}$.  In each
case the $b$ directions force the two parameter sets to be cosets of
a subgroup of order $b$: for finite sets $S,T$ in an abelian group,
$|S|=|T|=|S+T|$ implies $s+T=S+T$ for all $s\in S$, so $S$ and $T$
are cosets of the stabilizer of $T$.  The group orders are
$p,p-1,p+1$ respectively.  Since $1<b<p$, this proves the necessity.

The split and nonsplit constructions in Theorem~\ref{thm:digit} prove
sufficiency in every rank $r\ge2$.
\end{proof}

\subsection{A gap independent of rank}

\begin{lemma}[Separating a point in a nonplanar set]\label{base:separator}
Let $m\ge2$. For a set $S$ of at most $2m+2$ points in an affine space
over a field, assume $S$ is contained in no affine plane.
There are $P\in S$ and a polynomial $F$ of degree at most $m$ that
vanishes on $S\setminus\{P\}$ and is nonzero at $P$.
\end{lemma}

\begin{proof}
If a line contains at least $m+2$ points of $S$, choose $P$ outside it.
An affine hyperplane contains the line and avoids $P$. Each of the at most
$m-1$ other points outside the line and distinct from $P$ can be covered by a
hyperplane avoiding $P$.  Multiply the defining affine forms.

Otherwise fix any $P\in S$.  Group the nonzero vectors $x-P$ by their
one-dimensional spans.  Each group has at most $m$ members.  At most
two groups have size $m$, since there are at most $2m+1$ vectors and
$m\ge2$.  Choose a linearly independent triple meeting every group of
size $m$. This is possible since distinct one-dimensional subspaces
supply independent pairs, and the vectors together have rank at least
three.  The remaining vectors number at most $2(m-1)$ and each group now
has size at most $m-1$.  They partition into at most $m-1$ independent
pairs or singletons. Indeed, pad the list to $2(m-1)$ with dummy elements
in distinct new groups, sort by group, and pair entries $m-1$ places apart
across the two halves. Then discard the dummies.

For each resulting group, as well as for the triple, choose a linear form
taking value $1$ on its independent vectors.  Its affine level set
through the associated points avoids $P$.  The product of these at most
$m$ affine equations has the required properties.
\end{proof}

\begin{theorem}[The first gap below the character order]\label{base:gap}
Suppose $3\le b<p$, $r\ge2$, and $b$ divides neither $p-1$ nor $p+1$.
Every pair of distinct packets of size $b$ has disagreement density at
least $d_{b-1}=(p-1)(p+2-b)/p^2$.  This value is attained whenever
$b-1$ divides $p-1$ or $p+1$.
\end{theorem}

\begin{proof}
Write $S=\{x:c_x\ne0\}$ after cancellation, so $|S|\le2b$ and each
$c_x$ is nonzero in $\Fp$.  If $S$ is contained in an affine plane,
translate it into a vector plane. Translation adds a common constant
to both projected multisets and preserves agreement. Restricting the
uniform form $v$ to that plane gives a uniform plane form. The
disagreement cardinality there is divisible by $p-1$, since it is a
union of punctured lines and excludes the origin. Strict inequality
in Theorem~\ref{base:rigid} therefore gives the lower bound $d_{b-1}$.

If $S$ is nonplanar, Lemma~\ref{base:separator} with $m=b-1$ yields
$F$ for which $\sum_xc_xF(x)\ne0$ in $\Fp$. Since $\sum c_x=0$,
some monomial $x^\alpha$ of degree $1\le j\le b-1$ has a nonzero
signed sum. Form the homogeneous polynomial
\[
 M_j(v)=\sum_xc_x\,v(x)^j.
\]
The coefficient of $v^\alpha$ in dual coordinates is
$\binom{j}{\alpha}\sum_xc_xx^\alpha\ne0$ since $j<p$. Thus
$M_j$ is a nonzero polynomial function. Equal projected multisets
give $M_j(v)=0$. Applying Lemma~\ref{lem:projectivecount} to $M_j$,
with its scalar-invariant support, gives disagreement density
at least $d_j\ge d_{b-1}$.  Finally, when $b-1$ is an eligible
divisor, pad the extremal packets of size $b-1$ with a common entry.
\end{proof}

\begin{corollary}[An exact gap at character order seven]\label{base:seven}
For $p=7$, every $r\ge2$, and packets of size $5$, the minimum positive
disagreement density is $24/49$. Thus two distinct products may agree
at $25/49$ of the unramified primes, and no larger agreement density
is possible.
\end{corollary}

\begin{proof}
Five divides neither $6$ nor $8$, while four divides $8$.
Apply Theorem~\ref{base:gap}.
\end{proof}
\section{Equality at every scale}

The leading-digit bound has an exact criterion for equality. In particular,
increasing the rank cannot remove the obstruction from finite plane geometry.

\begin{theorem}[Equality and descent]\label{thm:descent}
Suppose $a\le r-2$. Then $D_p(r,k)$ is attainable if and only if
$b\mid(p-1)$ or $b\mid(p+1)$. In every equality case there is a subspace
$L\subset X$ of dimension $a$ such that, after cancellation of the common
entries, both packets are full inverse images under $X\to X/L$.
Their images are disjoint sets of size $b$ and lie in a common affine
plane. For $b\ge2$ they give a dual $3$-net as in Theorem~\ref{base:rigid}.
Conversely, every lift of a base equality case, padded equally on both
sides to size $k$, attains $D_p(r,k)$.
\end{theorem}

\begin{proof}
Let $d=a(p-1)+b$. The function $W(r-1,t-1)$ in
Lemma~\ref{lem:projectivecount} decreases strictly with the integer $t$
from $1$ to $d$ in this range. Hence equality forces $t=d$ and
$\operatorname{supp}H=\operatorname{supp}f$.

The reduction $f$ is represented by a homogeneous polynomial of degree $d$.
Indeed, for $u\in\Fp^\times$, Galois conjugation $\sigma_u(\zeta)=\zeta^u$
acts trivially on the residue field, and $\sigma_u(\pi)/\pi\equiv u\pmod\pi$.
Thus $f(uv)=u^d f(v)$. In the unique reduced representative with individual
exponents at most $p-1$, every monomial of nonzero coefficient consequently
has degree congruent to $d$ modulo $p-1$. Uniqueness follows by univariate
interpolation in successive variables. The constant term vanishes.
To homogenize each monomial, increase one of its positive exponents by
the difference between $d$ and its degree. This preserves its values
on $\Fp^r$, including at zeros of that coordinate.

Evaluation at projective representatives now gives a minimum-weight
word of the projective Reed--Muller code of degree $d$ on
$\mathbf P^{r-1}(\Fp)$. Its minimum weight is the value in
Lemma~\ref{lem:projectivecount}, attained here. By the classification
of Ghorpade and Ludhani \cite[Thm.~4.3]{GL}, this support is contained
in a linear space $U\subset V$ of codimension $a$ after taking its
affine cone. Explicitly their representing polynomial has a factor
\[
 L_a\prod_{i=0}^{a-1}(L_a^{p-1}-L_i^{p-1}),
\]
with independent linear forms. Wherever it is nonzero,
$L_a\ne0$ and $L_i=0$ for all $i<a$. Taking $U$ to be this simultaneous
kernel gives the assertion also when $a=0$.

Set $L=U^\perp$. Fourier inversion gives $c_{x+\ell}=c_x$ for every
$\ell\in L$, since the Fourier transform $H$ is supported in $U$.
Write $c'$ for the induced integer function on $X/L$. Its positive
and negative masses are equal and at most $\lfloor k/p^a\rfloor=b$.
On $U=(X/L)^*$ the transform $H$ is $p^a$ times the transform of $c'$,
and outside $U$ it vanishes. Thus the disagreement density for $c'$ is
$p^aD_p(r,k)=(p-1)(p+1-b)/p^2$. Theorem~\ref{base:rigid} applies
and gives all the stated restrictions. The lift construction of
Theorem~\ref{thm:digit} proves the converse.
\end{proof}

At powers $k=p^a$, before saturation, this has a particularly simple
interpretation. The extremizers are exactly two distinct parallel affine
$a$-spaces, with each character occurring once. At saturation a pair is
extremal exactly when its signed multiplicity function is constant on
cosets of some hyperplane in $X$ and is nonzero. Indeed, a minimal
disagreement set then consists of precisely one punctured line in $V$.
Inversion gives the stated periodicity, and conversely the support lies
on that line and contains all its nonzero points by Galois conjugacy.

\begin{corollary}[A gap at the smallest unsaturated rank]\label{cor:gaplift}
Let $r=a+2$. If $b$ divides neither $p-1$ nor $p+1$, and $b-1$ divides
one of them, the exact minimum is $(p-1)(p+2-b)/p^{a+2}$.
In particular, for $p=7$, $r=a+2$ and
$5\cdot 7^a\le k<6\cdot 7^a$, it equals $24/7^{a+2}$.
\end{corollary}

\begin{proof}
Density takes values in multiples of $(p-1)/p^r$.
Theorem~\ref{thm:descent} excludes the lower endpoint. The very next
multiple is attained by lifting the construction for the digit $b-1$
and padding with common entries.
\end{proof}
\section{Products whose factors all have the same primitive conductor}
\label{sec:primitive}

Requiring every factor to be primitive modulo $Q$ changes the
minimum density even when the packet size can grow arbitrarily.
Write $r$ for the dimension of the ambient exponent-$p$ character space
$X$ as before.

\begin{theorem}[Primitive common-conductor density floor]\label{thm:primitive}
If $A\ne C$ and every character in both packets is primitive modulo the same $Q$, then
\[
 \delta(A,C)\ \ge\ \frac{(p-1)2^{r-1}}{p^r}.
\]
For every modulus $Q$ supporting a primitive character of order $p$,
this is the exact minimum over packet sizes in the full space of
exponent-$p$ characters modulo $Q$, of rank $r=\omega(Q)$. It is
attained for every $k\ge2^{r-1}p^{r-1}$.
\end{theorem}

\begin{proof}
For odd $p$, the local characters of order dividing $p$ on a prime-power
factor of $Q$ form a group of dimension at most one: the units modulo
an odd prime power are cyclic, and those modulo a power of two form a
$2$-group. The local restrictions separate $X$, so choose $r$ independent
restrictions as coordinate functionals on $X$, choosing a generator of
each local order-$p$ group. A primitive character is nontrivial on every
local factor. Hence the signed multiplicity $c$ vanishes if any one of
its coordinates is zero. Fourier orthogonality now gives
\[
 \sum_{t\in\Fp} H(v+te_i)=0
 \qquad(1\le i\le r,\ v\in V)
\]
in dual coordinates.

A nonzero array on $\Fp^n$ with zero sums along every coordinate line
has at least $2^n$ nonzero entries. This follows by induction from the
case $n=0$: in any one coordinate at least two slices are nonzero by
the line-sum condition, and each such slice satisfies the remaining
conditions.  A nonzero slice of $H$ at $v_1\ne0$ exists by its first
line sums. Galois conjugation now makes all $p-1$ such slices nonzero.
Each contains at least $2^{r-1}$ nonzero entries.

For sharpness, such a modulus $Q$ is a product of distinct primes
$q\equiv1\pmod p$, possibly multiplied by $p^2$. Indeed for
$q\ne p$ an order-$p$ character factors through reduction modulo $q$,
whose kernel on units modulo any power of $q$ is a $q$-group. At $p$
the cyclic unit group modulo $p^e$ has its order-$p$ characters
already at $p^2$, and none at $p$. Consequently at every prime-power
factor of $Q$ the nontrivial local order-$p$ characters are exactly
the primitive ones of that order. Use these factors as coordinates and put
$\lambda(t)=1$ for even $r$ and
$\lambda(t)=\zeta^t-\zeta^{-t}$ for odd $r$, when $t\in\Fp^\times$.
Define $J(0,v')=0$ and
\[
 J(t,v')=\lambda(t)\prod_{i=2}^r
  \big(\mathbf1_{v_i=t}-\mathbf1_{v_i=-t}\big)
 \qquad(t\ne0).
\]
Each nonzero slice has $2^{r-1}$ nonzero entries.  Its line sums vanish in coordinates
$i\ge2$. In the first coordinate the only possible nonzero terms come
in a pair at $t,-t$ (when $r=1$ the terms cancel in opposite pairs).
The identity $\lambda(-t)=(-1)^r\lambda(t)$ gives cancellation.
Thus the inverse transform vanishes at characters having any zero
coordinate.  To see its integrality with an explicit size estimate,
define $c_x=p^{-1}\sum_{v\in V}J(v)\zeta^{-v(x)}$.
For $r\ge2$ expand the product over
$\epsilon\in\{1,-1\}^{r-1}$, write
$u_\epsilon=x_1+\sum_{i=2}^r\epsilon_i x_i$ and
$e_\epsilon=\prod_{i=2}^r\epsilon_i$. For even $r$,
\[
 c_x=\sum_\epsilon e_\epsilon\mathbf1_{u_\epsilon=0}.
\]
Indeed $\sum_{t\ne0}\zeta^{tu}=p\mathbf1_{u=0}-1$, and
$\sum e_\epsilon=0$. For odd $r$, including $r=1$ using the empty sign
vector, the formula is
\[
 c_x=\sum_\epsilon e_\epsilon
       (\mathbf1_{u_\epsilon=1}-\mathbf1_{u_\epsilon=-1}).
\]
Every indicator here is supported on a hyperplane of size $p^{r-1}$.
The coefficients are integers with sum zero, since $J(0)=0$, and
their Fourier transform is $p^{r-1}J$. Their positive and negative
masses are equal and at most $2^{r-1}p^{r-1}$. They therefore give
packets of the required size after common padding by a primitive
character. All nonzero coefficients have every coordinate nonzero, which
ensures primitivity at every prime-power factor of $Q$. Their disagreement set
has exactly $(p-1)2^{r-1}$ elements.
\end{proof}
\begin{corollary}[Shape of the primitive extremal support]\label{thm:primshape}
In the local coordinates used above, suppose the primitive bound is
attained. In rank two the disagreement set is the union of two
punctured lines of nonzero finite slopes.  In rank $r\ge3$ there are
$d_2,\dots,d_r\in\Fp^\times$ such that this set is exactly
\[
 \{(t,\epsilon_2d_2t,\dots,\epsilon_rd_rt):
      t\in\Fp^\times,\ \epsilon_i\in\{1,-1\}\}.
\]
\end{corollary}

\begin{proof}
A minimum-support array with all coordinate line sums zero in dimension
$n$ has support a Cartesian product of two-element sets. In the induction
above, equality forces exactly two nonzero slices. They are pointwise
opposites by the line-sum identity, so their supports agree and themselves
attain the minimum in dimension $n-1$. The assertion follows by induction.

Applying the primitive lower bound's counting argument with each coordinate
chosen first shows that the support avoids all coordinate hyperplanes.
The slice $v_1=1$ has support $D_2\times\cdots\times D_r$ with
each $D_i=\{a_i,b_i\}\subset\Fp^\times$. Every slice $v_1=t\ne0$
has support $tD_2\times\cdots\times tD_r$ by Galois conjugation. This
proves the rank-two description. If $r\ge3$, choose any $x_i\in D_i$.
The line with coordinates $v_i=x_i$ for $i\ge2$ must have a second
support point $t\ne1$ since its entries sum to zero. For every $i$,
$x_i/t$ then equals the other element of $D_i$. Consequently the ratios
of each chosen element to its partner are equal across all $i\ge2$.
Changing the choice in a single coordinate preserves another ratio and
inverts the changed ratio. The common ratio squares to one and differs
from one, so it is $-1$. Thus $D_i=\{d_i,-d_i\}$ for all $i$.
\end{proof}
\subsection{Exact densities for primitive cubic packets of any size}

The common-conductor result can be sharpened at $p=3$ for every $k$,
including sizes far below the saturation scale.

\begin{theorem}[Primitive cubic packets of prescribed size]\label{thm:cubic}
Let $p=3$, $k\ge1$, and $d=\min\{r,\lfloor\log_2(2k)\rfloor\}$.
Two distinct size-$k$ packets of primitive characters of one conductor
in a rank-$r$ space satisfy
\[
                  \delta(A,C)\ \ge\ (2/3)^d .
\]
For every conductor supporting a primitive cubic character, equality
is achievable in its full cubic character space, for every $k$.
In particular the first size at which the common-conductor minimum
$(2/3)^r$ is achieved is exactly $2^{r-1}$.
\end{theorem}

\begin{proof}
Use the independent local coordinates of Theorem~\ref{thm:primitive}.
Every nonzero signed multiplicity $c_x$ has $x_i\in\{1,2\}$.
Write $z_i=\zeta^{v_i}$. Dividing $H(v)$ by $\prod_i z_i$ gives a
multilinear polynomial in $z_1,\ldots,z_r$ on the cube of complex
cube roots of unity. It has $n\le2k$ monomials. We claim that a nonzero
multilinear polynomial with $n$ monomials on this grid is nonzero on a
fraction at least
\[
             (2/3)^{\min(r,\lfloor\log_2 n\rfloor)} .
\]
A single monomial never vanishes. If there is more than one monomial,
some coordinate has both exponents $0$ and $1$ in the support.
Write $P=A+z_i B$ in that coordinate. Both parts are nonzero, and
the smaller has at most $\lfloor n/2\rfloor$ monomials in $r-1$
variables. Wherever this part is nonzero, at most one of the three
choices for $z_i$ makes $P$ vanish. Applying the induction hypothesis
to the smaller part proves the claim. Increasing $n$ to $2k$ gives
the bound for the packets.

For equality in the full local character space, take the integer
coefficient array of
\[
    \left(\prod_{i=1}^r z_i\right)\prod_{i=1}^d(z_i-1).
\]
Its positive and negative coefficients give disjoint packets of
size $2^{d-1}$. All local exponents are $1$ or $2$, so all factors
are primitive. Its disagreement set consists of the vectors for
which $z_i\ne1$ for $1\le i\le d$, of proportion $(2/3)^d$.
Pad with a common primitive character to reach $k$.
\end{proof}

For example, in rank three the exact minimum disagreement densities for
primitive cubic packets of sizes $1,2,3,4$ are respectively
$2/3,4/9,4/9,8/27$. For unrestricted packets the same four values
are $2/3,4/9,2/9,2/9$. Thus three factors already distinguish the
two questions. For primitive packets the extremal construction is a
parity split of a Boolean cube of local character choices.
\section{Reading the distinction at finite primes}\label{sec:primes}

Fix a smooth nonnegative $w$ supported in $(1/2,1)$ with
$\int w=1$. For a function $F$ on $V$ and a real parameter $Y$, write
\[
 M_Y(F)=\sum_{\ell\nmid Q\ {\rm prime}}(\log\ell)w(\ell/Y)F(v_\ell),
 \qquad
 \|F\|_{A}=\sum_{x\in X}|\widehat F(x)|,
\]
where $v_\ell$ is evaluation at $\ell$ and the Fourier coefficients
use normalized counting measure on $V$. The article's prime Gram
estimate gives, using Dirichlet GRH, 
\begin{equation}\label{eq:prime}
 M_Y(F)=Y\,\mathbb E_V F+O_w\!\left(\|F\|_A\sqrt Y\log(2Q)\right).
\end{equation}
Indeed it suffices to apply the smoothed explicit formula to each
character of $V$, interpreted as a Dirichlet character modulo $Q$.
For its primitive inducing character, the weighted von Mangoldt sum
has the principal pole term, and the zero contribution is
$O_w(\sqrt Y\log(2Q))$ by rapid Mellin decay on the critical line and
the unit-interval zero count. Prime powers of exponent at least two
cost $O_w(\sqrt Y)$ by Chebyshev's bound. Deleting primes dividing $Q$
costs $O_w(\log(2Q))$. Summing the Fourier expansion gives
\eqref{eq:prime}. The constants below may depend on this fixed weight.

\begin{theorem}[Finite distinction and extremal counts]\label{thm:window}
For every distinct pair of size $k$, each interval $(Y/2,Y)$ contains a
distinguishing prime once $Y\ge C_w k^2\log^2(2Q)$, using the repository's
Dirichlet GRH input. If the pair attains $D_p(r,k)$ and $a\le r-2$,
write $S$ for its set of distinguishing evaluation vectors, $d'=p+1-b$,
and $\delta=D_p(r,k)$. Then
\begin{align}
 \|\mathbf1_S\|_A &= \frac{2b d'(p-1)}{p^2},\label{eq:exactnorm}\\
 M_Y(\mathbf1_S)&=Y\delta+
 O_w\!\left(2bp^a\delta\sqrt Y\log(2Q)\right).\label{eq:window}
\end{align}
For a pair attaining the saturated minimum the corresponding norm
is $2(1-1/p)$.
\end{theorem}

\begin{proof}
Apply \eqref{eq:prime} first to $F=|H|^2$.
Its mean is $\sum_x c_x^2$, and its Fourier norm is at most
$(\sum_x|c_x|)^2$ by convolution. Since the $c_x$ are integers,
$\sum c_x^2\ge\sum |c_x|$ and $\sum |c_x|\le2k$.
The weighted sum is therefore positive on the asserted scale.
Positivity forces one prime in the support of $w$ with $H(v_\ell)\ne0$.

For an extremal pair, Theorem~\ref{thm:descent} reduces $S$, inside
$U=L^\perp$, to the inverse image under a two-dimensional quotient
of a union of $d'$ punctured lines. This also covers $b=1$:
choose a second independent coordinate alongside the form for the
singleton difference. On $\Fp^2$ let $P$ denote that union.
The Fourier coefficient of $\mathbf1_P$ at the trivial character is
$d'(p-1)/p^2$. At a nontrivial character it is $(p-d')/p^2$ if that
character is trivial on a selected line, and $-d'/p^2$ otherwise.
The two multiplicities are $d'(p-1)$ and $b(p-1)$. Summing absolute
values gives \eqref{eq:exactnorm}. Pullback along a surjection preserves
the Fourier norm by orthogonality. Extension by zero from $U$ to $V$
also preserves it: every coefficient on $U$ gives $[V:U]$ equal
coefficients divided by $[V:U]$. Now
$\|\mathbf1_S\|_A/\delta=2bp^a$, proving \eqref{eq:window}.
At saturation $S$ is one punctured line and on the line the function
is $1-\mathbf1_{\{0\}}$, of norm $2(1-1/p)$.
\end{proof}

For fixed $p$ the absolute error in the measured proportion
$M_Y(\mathbf1_S)/Y$ is $O_{p,w}(\log(2Q)/\sqrt Y)$ for all equality
cases, independent of rank and packet size. If the packet size is unknown, given $a$ and the promise
of an unsaturated equality case, the digit $b$ can therefore be
recovered by rounding a single weighted prime count when
$Y\ge C_{p,w}p^{2a}\log^2(2Q)$: the possible densities for distinct
digits differ by at least $(p-1)/p^{a+2}$.
\section{The complete spectrum for two factors}

\begin{proposition}\label{prop:two}
At odd prime order $p$, the possible \emph{agreement} densities of two
distinct packets of size two, in ambient rank $r$, are as follows.
\[
\begin{array}{cl}
\toprule
 r & \text{distinct agreement densities} \\
\midrule
 1 & 1/p \\
 2 & 1/p,\quad (2p-1)/p^2,\quad 1/p^2 \\
 \ge3 & 1/p,\quad (2p-1)/p^2,\quad 1/p^2,\quad (2p-1)/p^3 \\
\bottomrule
\end{array}
\]
\end{proposition}

\begin{proof}
Write the packets as $\{s+u,s-u\}$ and $\{t+w,t-w\}$. Since $2$ is
invertible, their projections agree precisely on
\[
 (s-t)^\perp\cap\big((u-w)^\perp\cup(u+w)^\perp\big).
\]
Let $e=\dim\langle s-t\rangle\in\{0,1\}$. Inside the first hyperplane
(or the whole space when $e=0$), the conditional fraction is $1$
if either of the remaining forms vanishes identically, $1/p$ if their
nonzero restrictions are proportional, and $(2p-1)/p^2$ if they
are independent. Multiply by $p^{-e}$ and omit the fraction one
corresponding to identical packets. The three vectors
$s-t,u-w,u+w$ can be prescribed freely. Choosing their classes in
$X/\langle s-t\rangle$ with the indicated dependencies realizes
every entry allowed by its dimension.
\end{proof}

\section{Examples and reproducible checks}

\paragraph{Fifteen quintic factors.}
Take $Q=11\cdot31\cdot41=13981$. Let $\chi_1,\chi_2,\chi_3$ send the
least primitive roots $2,3,6$ in their respective prime factors to
$\exp(2\pi i/5)$, and inflate to $Q$. Write $(i,j,t)$ for
$\chi_1^i\chi_2^j\chi_3^t$. The two disjoint packets are
\[
 A=\{(1,0),(2,2),(2,3)\}\times\mathbf F_5,
 \qquad C=-A .
\]
Every factor has exact order five.
The first two coordinates describe the order-three subgroup on the
norm-one conic in $\mathbf F_5[\theta]/(\theta^2-2)$ and its negative.
The third coordinate runs freely. Of the $125$ evaluation vectors,
$12$ distinguish the packets. Each vector is represented by $96$
reduced residue classes modulo $Q$, giving $1152$ distinguishing
classes out of $12000$. The smallest distinguishing primes are
$3,79,163,173,191$. Through $100000$, precisely $906$ of the $9589$
unramified primes distinguish the packets. The density $12/125$ follows
from the exact residue-class count and the prime number theorem in
progressions. The finite prime count illustrates its meaning.

\paragraph{Five primitive factors of order seven.}
For $Q=29\cdot43$, choose a nontrivial order-seven character at each
prime and use them as the coordinate basis. In $\mathbf F_7^2$ take
\[
\begin{aligned}
 A&=\{(4,3),(2,3),(3,4),(3,2),(3,3)\},\\
 C&=\{(5,5),(1,5),(5,1),(1,1),(3,3)\}.
\end{aligned}
\]
All ten entries have nonzero exponents at both prime factors, so all
the characters are primitive of conductor $1247$. Translate by
$-(3,3)$ and cancel the common zero. The remaining sets are
$\{1,-1,i,-i\}$ and $(2+2i)\{1,-1,i,-i\}$ in
$\mathbf F_7[i]$, $i^2=-1$. The scalar $2+2i$ has norm one
and square $i$. Four of the eight projective directions agree,
and the origin always agrees. This gives agreement on $25/49$ of
the primes, the maximum by Corollary~\ref{base:seven}.

\paragraph{Exhaustive checks.}
The bundled Python program uses exact integer projection histograms.
It enumerates every packet in ten small parameter sets, giving full
disagreement distributions for $220\,813\,980$ unordered pairs in
aggregate. To count pairs it groups packet signatures on each subset
of projective directions. If $N_j$ pairs agree on exactly $j$ directions,
the sum of collision counts on the $i$-element subsets equals
$\sum_{j\ge i}\binom jiN_j$. Back substitution gives the full histogram.
Selected rows are shown below. All counts are in \texttt{checks.json}.
\begin{center}
\begin{tabular}{rrr r r c}
\toprule
$p$ & $r$ & $k$ & Packets & Unordered pairs & Minimum $\delta$\\
\midrule
3 & 2 & 4 & 495 & 122265 & $2/9$\\
3 & 3 & 2 & 378 & 71253 & $4/9$\\
3 & 3 & 3 & 3654 & 6674031 & $2/9$\\
5 & 2 & 3 & 2925 & 4276350 & $12/25$\\
5 & 2 & 4 & 20475 & 209602575 & $8/25$\\
\bottomrule
\end{tabular}
\end{center}

The program separately enumerates all primitive cubic packets of sizes
$1,2,3,4$ in ranks three and four, accounting for $7\,913\,653$ unordered
pairs across the eight settings. Each minimum agrees with
Theorem~\ref{thm:cubic}. For the signed construction in
Theorem~\ref{thm:primitive}, it checks the integer projections at every
evaluation vector for $p=3,5,7$ and $1\le r\le4$. The two explicit
Dirichlet examples are independently tested on all reduced residue
classes of their moduli.

\paragraph{Further questions.}
For an ineligible leading digit at higher rank and packet size,
the next attainable density can exceed the digit bound. The first
family left by the general results has $p=7$, $a\ge1$, $b=5$ and
$r>a+2$. Does the value $24/7^{a+2}$ persist throughout that family?
Even for $p=3,5$, the general second-smallest density and the full
range below twice the minimum remain to be determined. For character orders $p\ge5$, a sharp primitive bound for each individual packet size would refine Theorem~\ref{thm:primitive}. The exact
Fourier norm in Theorem~\ref{thm:window} suggests a related arithmetic
question: how large can the least distinguishing prime actually be
among extremal pairs of a given conductor and size?
\begin{thebibliography}{9}
\bibitem{Article} ProveIt Contributors,
\emph{Beyond the Riemann hypothesis for Dirichlet $L$-functions:
consequences, transfers and limits}, 26 November 2026.
See the subsection on matching Euler factors and the questions on pooled
zero statistics, and the prime Gram estimate in the same source file.
\href{https://github.com/VladimirReshetnikov/ProveIt/blob/2a9f8e3ee74a1824d57546f0b1b4f1cef5df93d6/Analysis/GeneralizedRiemannHypothesis/Article/sections/rec.tex}{Pinned recovery section}.
\href{https://github.com/VladimirReshetnikov/ProveIt/blob/2a9f8e3ee74a1824d57546f0b1b4f1cef5df93d6/Analysis/GeneralizedRiemannHypothesis/Article/sections/src.tex}{Source theorem and proof architecture}.

\bibitem{Lean} L.~Sernau and collaborators, readable Lean edition of the
Dirichlet GRH proof of L.~Alp\"oge, upstream commit
\href{https://github.com/lsernau/grh-lean/tree/72fead102f663d5c7a407964a13dc0cf9d977207/grh}{\texttt{72fead102f663d5c7a407964a13dc0cf9d977207}}.
ProveIt vendor and upstream share tree
\texttt{2e133df63d7f4da6b3059c4195324f8e9187ba73}.
\href{https://github.com/VladimirReshetnikov/ProveIt/blob/2a9f8e3ee74a1824d57546f0b1b4f1cef5df93d6/Analysis/GeneralizedRiemannHypothesis/vendor/grh-lean/grh/src/Statements.lean}{Definitions}.
\href{https://github.com/VladimirReshetnikov/ProveIt/blob/2a9f8e3ee74a1824d57546f0b1b4f1cef5df93d6/Analysis/GeneralizedRiemannHypothesis/vendor/grh-lean/grh/src/Main.lean}{Exported theorems}.
\href{https://github.com/VladimirReshetnikov/ProveIt/blob/2a9f8e3ee74a1824d57546f0b1b4f1cef5df93d6/Analysis/GeneralizedRiemannHypothesis/vendor/grh-lean/grh/README.md}{Build and verification report}.

\bibitem{Lev} V.~F.~Lev, Point distribution and perfect directions in
$\mathbf F_p^2$, \emph{Uniform Distribution Theory} \textbf{15} (2020),
93--98. \href{https://doi.org/10.2478/udt-2020-0012}{doi:10.2478/udt-2020-0012}.

\bibitem{BKM} A.~Blokhuis, G.~Korchm\'aros and F.~Mazzocca,
On the structure of $3$-nets embedded in a projective plane,
\emph{Journal of Combinatorial Theory, Series A} \textbf{118} (2011).
See Theorem~5.1 and its ensuing finite-field cases.
\href{https://arxiv.org/abs/0911.4100}{arXiv:0911.4100}.

\bibitem{GL} S.~R.~Ghorpade and R.~Ludhani, On the minimum distance,
minimum weight codewords, and the dimension of projective Reed--Muller
codes, \emph{Advances in Mathematics of Communications} \textbf{18}
(2024), 360--382. \href{https://doi.org/10.3934/amc.2023035}{doi:10.3934/amc.2023035}.

\bibitem{Chen} H.~Chen, On a generalization of Craig lattices,
\emph{Journal de Th\'eorie des Nombres de Bordeaux} \textbf{25} (2013),
59--70. \href{https://doi.org/10.5802/jtnb.825}{doi:10.5802/jtnb.825}.

\bibitem{Meshulam} R.~Meshulam, An uncertainty inequality for finite
abelian groups, \emph{European Journal of Combinatorics} \textbf{27}
(2006), 63--67. \href{https://arxiv.org/abs/math/0312407}{arXiv:math/0312407}.
\end{thebibliography}
\end{document}
